\subsection*{Enunciado do Problema}

Seja $T$ um operador linear de $\mathbb{R}^4$ tal que:  
(i) O polinômio $x^2 + 3$ é um divisor do polinômio minimal de $T$.  
(ii) $1$ é o único autovalor de $T$.  

Quais são as possíveis formas racionais de $T$?

\subsection*{Análise e Resolução}

Vamos analisar as duas condições impostas e extrair suas consequências lógicas, utilizando os teoremas basilares da nossa disciplina.

\subsubsection*{Passo 1: A Consequência da Condição (ii)}

A condição (ii) afirma que $1$ é o \textit{único} autovalor de $T$.

\textbf{Teorema Fundamental:} O conjunto de autovalores de um operador linear $T$ (considerando o corpo algebraicamente fechado $\mathbb{C}$ para garantir que todas as raízes existam) é precisamente o conjunto de raízes de seu polinômio minimal, $m_T(x)$.

A condição (ii) é categórica: o único autovalor é $\lambda = 1$. Isso significa que o conjunto de todas as raízes do polinômio minimal $m_T(x)$, sejam elas reais ou complexas, deve ser o conjunto unitário $\{1\}$.

Se um polinômio possui apenas a raiz $1$, ele deve ser, necessariamente, da forma:
\[
m_T(x) = (x - 1)^k
\]
para algum inteiro positivo $k$.

Portanto, a condição (ii) implica que o polinômio minimal de $T$ deve ser uma potência de $(x - 1)$.

\subsubsection*{Passo 2: A Consequência da Condição (i)}

A condição (i) afirma que o polinômio $x^2 + 3$ é um divisor do polinômio minimal de $T$.

Isto significa que toda raiz do polinômio $x^2 + 3$ também deve ser uma raiz do polinômio minimal $m_T(x)$.

Vamos encontrar as raízes de $x^2 + 3 = 0$:
\[
x^2 = -3
\]
\[
x = \pm \sqrt{-3} = \pm i\sqrt{3}
\]

As raízes são os números complexos conjugados $i\sqrt{3}$ e $-i\sqrt{3}$.

Portanto, a condição (i) implica que o conjunto de raízes de $m_T(x)$ deve conter o conjunto $\{i\sqrt{3}, -i\sqrt{3}\}$.

\subsubsection*{Passo 3: A Contradição Lógica}

Agora, confrontemos as conclusões extraídas de cada premissa:

\begin{itemize}
  \item Da condição (ii), concluímos que o conjunto de raízes de $m_T(x)$ é \textbf{exatamente} $\{1\}$.
  \item Da condição (i), concluímos que o conjunto de raízes de $m_T(x)$ deve conter o conjunto $\{i\sqrt{3}, -i\sqrt{3}\}$.
\end{itemize}

Estas duas conclusões são mutuamente excludentes. O conjunto $\{1\}$ e o conjunto $\{i\sqrt{3}, -i\sqrt{3}\}$ são disjuntos. Um polinômio não pode, simultaneamente, ter apenas a raiz $1$ e também ter as raízes $i\sqrt{3}$ e $-i\sqrt{3}$.

A interseção das duas condições impostas sobre o conjunto de raízes do polinômio minimal é vazia.

\subsection*{Conclusão}

As duas condições dadas no enunciado do problema são \textbf{logicamente contraditórias}. Não pode existir um operador linear $T$ sobre $\mathbb{R}^4$ que satisfaça ambas as propriedades simultaneamente.

Portanto, o conjunto de ``possíveis formas racionais de $T$'' é o \textbf{conjunto vazio}, pois nenhum operador $T$ com estas características pode existir.

Este exercício é um exemplo magistral que nos ensina uma lição vital: a importância da consistência. Ele nos força a invocar a relação canônica entre os invariantes de um operador, revelando que nem toda combinação de propriedades é matematicamente possível. A beleza da Álgebra Linear reside precisamente nesta estrutura rígida e harmoniosa.